\chapter{积分变换法}
\begin{definition}{积分变换}{Integral Transformation}
	积分变换
	\begin{equation}
		\label{eq:mmp-integral-transform}
		\mathcal T\bigb{f}(\xi) = \int_a^bf(x)K(x,\xi)\d x.
	\end{equation}
	其中$K$是核函数(kernel)。
\end{definition}
\section{Fourier变换}
\setcounter{subsubsection}{1}
\begin{definition}{Fourier变换}{Fourier Transformation}
	定义Fourier变换
	\begin{equation}
		\label{eq:mmp-fourier-transform}
		\hat f(\xi) \equiv \FT[f](\xi) := \intR f(x)\e^{-\i\xi x}\d x,\quad \xi \in \RR,
	\end{equation}
	及其逆变换
	\begin{equation}
		\label{eq:mmp-inverse-fourier-transform}
		f(x) = \FT\inv\big[\hat f\big](x) = \frac1{2\pi}\intR\hat f(\xi)\e^{\i\xi x}\d\xi.
	\end{equation}
\end{definition}
Fourier变换的相关性质见第 \pageref{The Property of Fourier Transformation} 页的附录内容。
\begin{example}{无限长杆热传导}{}
	无限长杆热传导
	\begin{align*}
		\begin{cases}
			u_t = a^2u_{xx} + f(x,t),\\
			u(x,0) = \varphi(x).
		\end{cases} \longrightarrow \quad\begin{cases}
			\hat u_t = -a^2\xi^2\hat u + \hat f,\\
			\hat u(\xi,0) = \hat\varphi(\xi).
		\end{cases}
	\end{align*}
	因此
	\begin{align*}
		\hat u(\xi,t) &= \e^{-a^2\xi^2t}\biggb{\int_0^t\hat f(\xi,\tau)\e^{a^2\xi^2\tau}\d\tau + \hat\varphi(\xi)}.
		% \\
		% &=\int_0^t\hat f(\xi,\tau)\e^{-a^2\xi^2(t-\tau)}\d\tau+\hat\varphi(\xi)\e^{-a^2\xi^2t}.
	\end{align*}
	再用Fourier逆变换
	\begin{align}
		u(x,t) &= \FT\inv\biggb{\int_0^t\hat f(\xi,\tau)\e^{-a^2\xi^2(t-\tau)}\d\tau} + \FT\inv\bigb{\hat\varphi(\xi)\e^{-a^2\xi^2t}}\notag\\
		&= \int_0^t\FT\inv\bigb{\hat f(\xi,\tau)} \ast \FT\inv\bigb{\e^{-a^2\xi^2(t-\tau)}}\d\tau + \FT\inv\bigb{\hat\varphi(\xi)} \ast \FT\inv\bigb{\e^{-a^2\xi^2t}}\notag\\
		&= \int_0^tf(x,\tau) \ast \frac{\e^{-x^2/4a^2(t-\tau)}}{2a\sqrt{\pi(t - \tau)}}\d\tau + \frac1{2a\sqrt{\pi t}}\intR\varphi(\eta)\e^{-(x-\eta)^2/4a^2t}\d\eta\notag\\
		&= \int_0^t\frac1{2a\sqrt{\pi(t - \tau)}}\intR f(\eta,\tau)\e^{-(x-\eta)^2/4a^2(t-\tau)}\d\eta\d\tau +\notag\\
		&\qqquad\frac1{2a\sqrt{\pi t}}\intR\varphi(\eta)\e^{-(x-\eta)^2/4a^2t}\d\eta.\label{eq:mmp-forced-heat-kernel-solution}
	\end{align}
	以上只是形式解。
\end{example}
\subsection{Fourier正余弦变换}
对于定义在$\RR_{\geqslant  0}$上的函数，可将其奇偶延拓至$\RR$后再作Fourier变换
\begin{definition}{Fourier正余弦变换}{Fourier Sine and Cosine Transformation}
	正余弦变换分别为
	\begin{subequations}
		\label{eq:mmp-sine-cosine-transforms}
		\begin{align}
			\hat f_\sine(\xi) \equiv \FT_\sine[f](\xi) := \intRplus f(x)\sin\xi x\d x;\\
			\hat f_\cosine(\xi) \equiv \FT_\cosine[f](\xi) := \intRplus f(x)\cos\xi x\d x.
		\end{align}
	\end{subequations}
	其反变换
	\begin{equation}
		\label{eq:mmp-inverse-half-range-transforms}
		\begin{aligned}
			f(x) &= \frac2\pi\intRplus\hat f_\sine(\xi)\sin\xi x\d\xi =: \FT_\sine\inv[\hat f_\sine(\xi)]\\
			&= \frac2\pi\intRplus\hat f_\cosine(\xi)\cos\xi x\d\xi =: \FT_\cosine\inv[\hat f_\cosine(\xi)]
		\end{aligned}
	\end{equation}
\end{definition}
变换的性质见第 \pageref{The Property of Sine and Cosine Transformation} 页的附录内容。
\begin{example}{半无限长杆热传导}{}
	半无限长杆热传导方程问题
	\begin{align*}
		\begin{cases}
			u_t = a^2u_{xx},\quad x > 0,\;t > 0\\
			u(x,0) = 0,\\
			u_x(0,t) = \varphi(t)\\
		\end{cases}
	\end{align*}
	自然边界条件
	\[
		\lim_{x\to\infty}u(x,t) = \lim_{x\to\infty}u_x(x,t) = 0
	\]
	
	若用Fourier正弦变换，则会出现$u(0,t)$，此边界条件没有给出，因此只能用余弦变换，记$U(\xi,t) = \FT_\cosine[u(x,t)]$
	\begin{align*}
		U_t(\xi,t) &= a^2\bigb{u_x(x,t)\cos\xi x + \xi u(x,t)\sin\xi x}_0^\infty - a^2\xi^2U(\xi,t)\\
		&= -a^2\varphi(t) - a^2\xi^2U(\xi,t).
	\end{align*}
	故
	\[
		U(\xi,t) = -a^2\int_0^t\varphi(\tau)\e^{-a^2\xi^2(t-\tau)}\d\tau
	\]
	逆变换
	\begin{align}
		u(x,t) &= -\frac{2a^2}\pi\intRplus\biggb{\int_0^t\varphi(\tau)\e^{-a^2\xi^2(t-\tau)}\d\tau}\cos\xi x\d\xi\notag\\
		&= -\frac{2a^2}\pi\int_0^t\varphi(\tau)\biggp{\intRplus\e^{-a^2\xi^2(t-\tau)}\cos\xi x\d\xi}\d\tau\notag\\
		&= -a\int_0^t\varphi(\tau)\frac{\e^{-x^2/4a^2(t-\tau)}}{\sqrt{\pi(t - \tau)}}\d\tau.\label{eq:mmp-half-line-heat-flux}
	\end{align}
\end{example}
\begin{example}{半无限长弦振动}{}
	半无界弦振动问题
	\begin{align*}
		\begin{cases}
			u_{tt} = a^2u_{xx},\quad x > 0,\;t > 0\\
			u_x(0,t) = f(t),\\ % \quad u(0,t)=f(t),
			u(x,0) = u_t(x,0) = 0,\\
			\lim_{x\to\infty}u(x,t) = \lim_{x\to\infty}u_x(x,t) = 0.
		\end{cases}
	\end{align*}
	最后一行是无穷远处的自然边界条件。对空间坐标作Fourier余弦变换，得到
	\[
		U_{tt}(\xi,t) = -a^2f(t) - a^2\xi^2U(\xi,t).
	\]
	齐次方程通解$\cos a\xi t,\;\sin a\xi t$，运用\eqref{General to Particular Solution}和边界条件$U(\xi,0) = U_t(\xi,0) = 0$
	\begin{align*}
		U(\xi,t) &= \int_0^t - a^2f(\tau)\frac{\cos a\xi\tau\sin a\xi t - \cos a\xi t\sin a\xi\tau}{\cos a\xi\tau(\sin a\xi\tau)' - (\cos a\xi\tau)'\sin a\xi\tau}\d\tau\\
		&= -\frac a\xi\int_0^tf(\tau)\sin a\xi(t - \tau)\d\tau.
	\end{align*}
	逆变换
	\begin{align*}
		u(x,t) &= -\frac{2}\pi\intRplus\biggb{\frac a\xi\int_0^tf(\tau)\sin a\xi(t - \tau)\d\tau}\cos \xi x\d\xi\\
		&= -2a\int_0^tf(\tau) \cdot \frac1{2\pi}\intR\frac1\xi\sin a\xi(t - \tau)\e^{\i\xi x}\d\xi\d\tau
	\end{align*}
	由
	\[
		\FT\inv\Bigb{\frac{\sin a\xi}\xi}(x) = \frac{\Heaviside(x + a) - \Heaviside(x - a)}2,\quad a > 0.
	\]
	其中$\Heaviside(x)$为Heaviside跳跃函数，故
	\begin{equation}
		\label{eq:mmp-half-line-wave-flux}
		u(x,t)  %=-a\int_0^tf(\tau)[H(x+a(t-\tau))-H(x-a(t-\tau))]\d\tau\\
		% =-a\int_0^{t-x/a}f(\tau)\Heaviside(x+a(t-\tau))\d\tau
		= \begin{cases}-a\int_0^{t-x/a}f(\tau)\d\tau,&0 \leq x < at,\\0,&x \geq at.\end{cases}
	\end{equation}
对足够光滑的$f$，若要求跨越波前$x = at$的古典解，还需初边值相容条件$f(0) = f'(0) = 0$；否则应在弱解意义下理解波前处的导数。
\end{example}
\subsection{Fourier变换与分离变量法}
无界区域是否有相应的分离变量法？
\begin{example}{无界长杆热传导}{}
	
	无界长杆热传导
	\begin{align*}
		\begin{cases}
			u_t = a^2u_{xx},\quad x \in \RR,\;t > 0\\
			u(x,0) = \varphi(x)
		\end{cases}
	\end{align*}
	分离变量$u(x,t) = X(x)T(t)$特征值
	\[
		X'' + \lambda X = 0,\;T' + \lambda a^2T = 0
	\]
	讨论知$\lambda = \xi^2 \geqslant 0$
	\[
		X = \e^{\i\xi x},\quad T = \e^{-a^2\xi^2t}.
	\]
	故
	\begin{align*}
		u(x,t) &= \frac1{2\pi}\intR A(\xi)X(x,\xi)T(t,\xi)\d\xi\\
		&= \frac1{2\pi}\intR A(\xi)\e^{\i\xi x}\e^{-a^2\xi^2t}\d\xi.
	\end{align*}
	边界条件
	\[
		u(x,0) = \frac1{2\pi}\intR A(\xi)\e^{\i\xi x}\d\xi \equiv \FT\inv[A] = \varphi(x),
	\]
	故$\varphi(x)$是$A(\xi)$的Fourier逆变换，即$A = \FT[\varphi]$。
\end{example}
可见对无界区间上的问题，分离变量法依然适用，只是特征值是连续分布的。
\section{Laplace变换}
\begin{definition}{Laplace变换}{Laplace Transformation}
	定义Laplace变换
	\begin{equation}
		\label{eq:mmp-laplace-transform}
		\bar f(\xi) \equiv \LT[f](\xi) := \intRplus f(x)\e^{-\xi x}\d x,\quad\Re\xi > \sigma_0.
	\end{equation}
	
	其逆变换
	\begin{equation}
		\label{eq:mmp-bromwich-inversion}
		f(x) = \LT\inv[\bar f(\xi)] = \frac1{2\pi\i}\int_{\sigma-\i\infty}^{\sigma+\i\infty}\bar f(\xi)\e^{\xi x}\d\xi,\quad x \geqslant 0.
	\end{equation}
\end{definition}
Laplace变换存在的一组充分条件是$f$局部可积，且$\exists M,\sigma_0$使得
	\[
		\abs{f(x)} < M\e^{\sigma_0x},\quad\forall x > 0.
	\]
	% $x$为$f$连续点处
% 再看Laplace逆变换$\LT\inv[\bar f(\xi)]\to f(x)$。
% \begin{theorem}{Riemann-Mellin反演公式}{Riemann-Mellin Inversion Formula}
其逆变换中$\sigma > \sigma_0$；注意，Laplace变换的积分中本身不包含$\RR_{<0}$的部分，一般认为$f(< 0) \equiv 0$。
\begin{theorem}{第一展开定理}{}
	设$F(\xi)$在$\infty$邻域内有Laurent展开式
	\[
		F(\xi) = \sum_{n=1}^\infty\frac{c_n}{\xi^n},
	\]
	则
	\begin{equation}
		\label{eq:mmp-laplace-series-inverse}
		f(x) = \LT\inv[F(\xi)] = \sum_{n=1}^\infty\frac{c_n}{(n - 1)!}x^{n-1},\quad x \geqslant 0
	\end{equation}
\end{theorem}
\begin{theorem}{第二展开定理}{}
	设$F(\xi) = A(\xi)/B(\xi)$是有理函数，$\deg A < \deg B$，$B(\xi)$只有单零点$\xi_1,\xi_2,\ldots,\xi_n$且是$F(\xi)$的单极点，则
	\begin{equation}
		\label{eq:mmp-laplace-simple-poles}
		f(x) = \LT\inv[F(\xi)] = \sum_{k=1}^n\frac{A(\xi_k)}{B'(\xi_k)}\e^{\xi_kx} = \sum_{k=1}^n\Res\bigb{F(\xi)\e^{\xi x},\xi_k}.
	\end{equation}
\tcblower
	若$F(\xi)$满足
	\begin{compactenum}
		\item 在$\CC$上除了有限个奇点$\xi_1,\xi_2,\ldots,\xi_n$外解析；
		\item 在半平面$\Re\xi > \sigma_0$上解析；
		\item $\exists M > 0,\;R > 0$使得当$\abs\xi > R$时
		\[
			\abs{F(\xi)} \leqslant \frac M{\abs\xi},
		\]
	\end{compactenum}
	则对于$x \geqslant 0$有
	\begin{equation}
		\label{eq:mmp-laplace-residue-inversion}
		f(x) = \LT\inv[F(\xi)] = \sum_{k=1}^n\Res[F(\xi)\e^{\xi x},\xi_k].
	\end{equation}
\end{theorem}
% \subsection{Laplace定理的应用}
\begin{example}{有限杆热传导}{}
	有限杆热传导
	\[
		\begin{cases}
			u_t = u_{xx},\quad x \in (0,L),\;t > 0,     \\
			u_x(0,t) = 0,\quad u(L,t) = A, \\
			u(x,0) = B.
		\end{cases}
	\]
	关于$t$作Laplace变换，记$U(x,s) = \LT[u(x,t)](s)$
	\[
		s U(x,s) - u(x,0) = U_{xx}(x,s),
	\]
	由$u(x,0) = B$，
	\[
		U(x,s) = c_1\cosh\sqrt s x + c_2\sinh\sqrt s x + \frac Bs.
	\]
	再由$U_x(0,s) = 0,\;U(L,s) = A/s$
	\[
		U(x,s) = \frac{(A - B)\cosh\sqrt s x}{s\cosh\sqrt s L} + \frac Bs.
	\]
	逆变换
	\begin{align*}
		u(x,t) &= \LT\inv\biggb{\frac{(A - B)\cosh\sqrt s x}{s\cosh\sqrt s L}}(t) + \LT\inv\biggb{\frac Bs}(t)\\
		&= (A - B)\LT\inv\biggb{\frac{\cosh\sqrt s x}{s\cosh\sqrt s L}}(t) + B.
	\end{align*}
	括号内函数的孤立奇点\footnote{注意，0不是此函数的孤立奇点，不能用留数算出。逆变换结果中的1是围道积出来的。详情可见附录}应使得$\cosh\sqrt s L = 0$
	\[
		s_n = -\biggb{\frac{(2n - 1)\pi}{2L}}^2,\quad n = 1,2,\ldots
	\]
	由
	\begin{align}
		u(x,t) &= B + (A - B)\biggp{1 + \sum_{n=1}^\infty\Res\biggb{\frac{\cosh\sqrt s x}{s\cosh\sqrt s L}\e^{s t},s_n}}\notag\\
		&= A + \frac{4(A - B)}\pi\sum_{n=1}^\infty\frac{(-1)^n}{2n - 1}\cos\biggb{\frac{(2n - 1)\pi}{2L}x}\exp\biggb{-\frac{(2n - 1)^2\pi^2}{4L^2}t}.\label{eq:mmp-mixed-rod-temperature}
	\end{align}
\end{example}
\begin{example}{有限长弦受迫振动}{}
	有限长弦受迫振动
	\begin{align*}
		\begin{cases}
			u_{tt} = a^2u_{xx},\quad x \in (0,L),\;t > 0\\
			u(0,t) = 0,\quad u_t(L,t) = A\sin\omega t\\
			u(x,0) = u_t(x,0) = 0
		\end{cases}
	\end{align*}
	关于$t$作Laplace变换，记$U(x,s) = \LT[u(x,t)](s)$
	\[
		s^2U(x,s) - s u(x,0) - u_t(x,0) = a^2U_{xx}(x,s).
	\]
	结合$U_x(L,s) = A\omega/(s^2 + \omega^2)$
	\[
		U(x,s) = \frac{Aa\omega}{s(s^2 + \omega^2)}\sinh\frac{xs}a\sech\frac{Ls}a.
	\]
	其所有的孤立奇点：可去奇点0；一级奇点$\pm\i\omega_k$，
	\[
		\omega_k :=
		\begin{cases}
			\omega,&k = 0\\
			\frac{2k - 1}{2L}\pi a,&k \geqslant 1
		\end{cases}
	\]
	故
	\begin{align}
		&u(x,t) = \sum_{k=0}^\infty\Res[U(x,s)\e^{st},\i\omega_k] + \Res[U(x,s)\e^{st},-\i\omega_k]\notag\\
		&= 2\Re\sum_{k=0}^\infty\Res[U(x,s)\e^{st},\i\omega_k]\notag\\
		&= 2\Re\Bigb{\frac{Aa\omega}{s(s + \i\omega)}\sinh\frac{xs}a\sech\frac{Ls}a\e^{st}}_{\i\omega} +\notag\\
		&\qquad 2\Re\sum_{k=1}^\infty\Bigb{\frac{Aa\omega}{s(s^2 + \omega^2)}\sinh\frac{xs}a \cdot \frac aL\csch\frac{Ls}a\e^{st}}_{\i\omega_k}\notag\\
		&= \frac{Aa}\omega\sinh\frac{\omega x}a\sec\frac{\omega L}a\sin\omega t +\notag\\
		&\qquad 16Aa\omega^2L\sum_{k=1}^\infty\frac{(-1)^{k-1}}{(2k - 1)[4\omega^2L^2 - (2k - 1)^2\pi^2a^2]}\sin\frac{\omega_kx}a\sin\omega_kt.\label{eq:mmp-driven-string-solution}
	\end{align}
\end{example}

